\subsection{Xây dựng vùng giao thông động}
\label{subsec:dynamic_zone_construction}

Như đã trình bày ở Mục~\ref{subsec:problem_overview}, vùng giao thông trong nghiên cứu này
không được xem là một khu vực không gian cố định, mà là một \textit{vùng chức năng động}
(functional dynamic zone) được khởi tạo từ các sự kiện ùn tắc và mở rộng dựa trên mối liên hệ
không gian--thời gian giữa các tuyến đường. Quá trình xây dựng vùng giao thông động bao gồm
hai giai đoạn chính: (i) phát hiện cụm ùn tắc ban đầu tại từng thời điểm, và (ii) mở rộng vùng
dựa trên cấu trúc mạng lưới và các tiêu chí không gian.

\subsubsection{Phát hiện cụm ùn tắc theo snapshot}
\label{subsubsec:dbscan_snapshot}

Tại mỗi thời điểm $t$, nhóm xác định tập các tuyến đường đang ở trạng thái ùn tắc dựa trên
nhãn nhị phân $y_i(t) \in \{0,1\}$. Đối với mỗi tuyến đường, tọa độ không gian được biểu diễn
bằng điểm trung tâm của tuyến, được tính từ trung điểm giữa hai nút đầu và cuối của tuyến
trong đồ thị giao thông.

Trên tập các tuyến đang ùn tắc tại thời điểm $t$, nhóm áp dụng thuật toán phân cụm
DBSCAN để phát hiện các \textit{cụm ùn tắc theo snapshot}. DBSCAN được lựa chọn do không
yêu cầu số cụm cố định trước và có khả năng phát hiện các cụm với hình dạng bất kỳ, phù hợp
với đặc điểm phân bố không gian của ùn tắc giao thông.
Mỗi cụm thu được từ DBSCAN tại thời điểm $t$ được xem như một \textit{seed congestion cluster},
đóng vai trò khởi tạo cho quá trình xây dựng vùng giao thông động.\
\subsubsection{Chiến lược chọn seed vùng}
\label{subsubsec:seed_sampling}

Để đảm bảo tính cân bằng giữa các trạng thái giao thông khác nhau trong quá trình huấn luyện,
nhóm áp dụng chiến lược chọn seed vùng kết hợp giữa các cụm ùn tắc và các vùng không
ùn tắc. Cụ thể, trong quá trình sinh mẫu huấn luyện, $60\%$ seed được lấy từ các cụm ùn tắc
được phát hiện bởi DBSCAN tại thời điểm $t$, trong khi $40\%$ seed còn lại được chọn từ các
vùng không ùn tắc.

Các seed không ùn tắc có thể được tạo bằng cách áp dụng DBSCAN trên tập các tuyến không
ùn tắc. Sau khi đã chọn ra các cụm kẹt hoặc không kẹt từ DBSCAN, mô hình sẽ lựa chọn ngẫu nhiên một vùng cụ thể với điều kiện
tương ứng tại thời điểm t. Chiến lược này giúp mô hình học được cả các
mẫu lan truyền ùn tắc lẫn các trạng thái giao thông ổn định, đồng thời giảm thiểu hiện tượng
lệch lớp trong dữ liệu.

\subsubsection{Mở rộng vùng theo cấu trúc mạng lưới}
\label{subsubsec:graph_expansion}

Từ mỗi seed congestion cluster, nhóm tiến hành mở rộng vùng giao thông để thu thập
tập các tuyến ứng viên có khả năng liên quan đến động lực giao thông của vùng. Trước hết,
tất cả các tuyến thuộc seed cluster được đưa vào tập vùng ban đầu. Sau đó, dựa trên cấu trúc
đồ thị giao thông toàn cục, vùng được mở rộng bằng cách bổ sung các tuyến lân cận theo quan
hệ topology, với số bước mở rộng (hop) được được nhóm chọn là 1.

Quá trình mở rộng này nhằm đảm bảo rằng vùng giao thông không chỉ bao gồm các tuyến đang
ùn tắc tại thời điểm hiện tại, mà còn bao phủ các tuyến có khả năng đóng vai trò trung gian
hoặc chịu ảnh hưởng từ quá trình lan truyền ùn tắc trong tương lai.

\subsubsection{Lọc vùng theo bán kính không gian}
\label{subsubsec:radius_filter}

Để giới hạn phạm vi không gian của vùng giao thông và đảm bảo tính hợp lý về mặt vật lý,
nhóm áp dụng một bước lọc cứng dựa trên bán kính không gian. Cụ thể, đối với mỗi seed
cluster, một điểm trung tâm của vùng được xác định từ trung tâm hình học của các tuyến
trong seed.

Với mỗi tuyến ứng viên thu được từ bước mở rộng theo topology, khoảng cách Haversine từ
tuyến đó đến tâm seed được tính toán. Chỉ những tuyến có khoảng cách nhỏ hơn hoặc bằng
ngưỡng $R_{\max}$ mới được giữ lại trong vùng giao thông cuối cùng. Ở mô hình này nhóm chọn bán kính là 1500m cho mỗi vùng. Bước lọc này giúp loại
bỏ các tuyến ở quá xa seed cluster, vốn khó có mối liên hệ trực tiếp với sự kiện ùn tắc ban đầu.

\subsubsection{Vùng giao thông chức năng}
\label{subsubsec:functional_zone}

Sau các bước trên, vùng giao thông thu được không còn là một cụm không gian thuần túy,
mà là một \textit{vùng giao thông chức năng}, bao gồm các tuyến đường có mối liên hệ về
không gian, cấu trúc mạng lưới và tiềm năng ảnh hưởng giao thông. Đáng chú ý, các tuyến
nằm trong vùng giao thông này không nhất thiết phải thuộc seed congestion cluster ban đầu,
mà có thể nằm ngoài cụm nhưng vẫn có vai trò quan trọng trong động lực lan truyền ùn tắc.

Vùng giao thông chức năng này đóng vai trò là đơn vị đầu vào cho các bước tiếp theo, trong
đó các mối tương quan giao thông có độ trễ sẽ được khai thác để xác định các liên kết chức
năng và tập các tuyến mục tiêu cần dự đoán.
